\section{Transacciones Distribuidas}

Una compra toca dos microservicios: \texttt{ServicioPagos} debita la plata y
\texttt{ServicioInventario} asigna el asiento. Son procesos distintos y cada
uno tiene su base. No hay una transacción local que abarque a los dos, así
que la atomicidad hay que armarla con mensajes. El protocolo es
\emph{Two-Phase Commit}. El código está en el módulo \texttt{dos\_fases} del
crate \texttt{tp2\_reservas} y esta parte se corre con
\texttt{cargo run -- parte3}.

Los dos servicios están armados como actores: un hilo con el recurso adentro
y un buzón \texttt{mpsc}. La diferencia está en quién decide. Un actor que
atiende un débito contra una sola cuenta resuelve ahí mismo si entra o no.
Acá su respuesta pasa a ser un voto y la decisión la toma el coordinador
cuando tiene los votos de todos.

\begin{figure}[H]
\centering
\begin{tikzpicture}[x=3.4cm,y=0.8cm]
  % Lineas de vida: el coordinador va en el medio para que ninguna flecha cruce
  \node[caja] (pagos) at (0,0) {ServicioPagos};
  \node[caja] (coord) at (1,0) {Coordinador};
  \node[caja] (inv) at (2,0) {ServicioInventario};
  \foreach \x in {0,1,2} \draw[dashed] (\x,-0.4) -- (\x,-5.6);

  % Fase 1: Prepare y votacion
  \draw[flecha] (1,-1.1) -- node[above,font=\scriptsize]{\texttt{Prepare}} (0,-1.1);
  \draw[flecha] (1,-1.1) -- node[above,font=\scriptsize]{\texttt{Prepare}} (2,-1.1);
  \draw[flecha] (0,-2.3) -- node[above,font=\scriptsize]{\texttt{Yes}} (1,-2.3);
  \draw[flecha] (2,-2.3) -- node[above,font=\scriptsize]{\texttt{Yes}} (1,-2.3);

  % Fase 2: decision
  \draw[flecha] (1,-3.9) -- node[above,font=\scriptsize]{\texttt{Commit}} (0,-3.9);
  \draw[flecha] (1,-3.9) -- node[above,font=\scriptsize]{\texttt{Commit}} (2,-3.9);

  % Etiquetas al costado
  \node[font=\scriptsize,anchor=west] at (2.25,-1.7) {FASE 1: votación};
  \node[font=\scriptsize,anchor=west] at (2.25,-3.1) {punto de decisión};
  \node[font=\scriptsize,anchor=west] at (2.25,-3.9) {FASE 2: commit};
  \node[font=\scriptsize,anchor=west] at (2.25,-4.9) {fondos y asiento firmes};
  \node[font=\scriptsize,anchor=east] at (-0.1,-3.1) {PREPARADO};
\end{tikzpicture}
\caption{Transacción exitosa: los dos participantes votan \texttt{Yes},
quedan PREPARADOS y esperan la decisión del coordinador.}
\label{fig:2pc-secuencia}
\end{figure}

\subsection{Los participantes y el estado PREPARADO}

Cada participante maneja su recurso con tres operaciones.
\texttt{preparar} reserva de forma tentativa y devuelve el voto,
\texttt{confirmar} hace definitiva la reserva y \texttt{revertir} la deshace.
En \texttt{ServicioPagos} preparar descuenta el monto y lo marca retenido,
confirmar levanta la retención y revertir devuelve la plata. En
\texttt{ServicioInventario} preparar marca el asiento como no libre y
retenido, confirmar lo pasa a vendido y revertir lo libera.

El voto \texttt{Yes} y la reserva son la misma operación. Un participante que
vota \texttt{Yes} promete que puede cumplir, y si el asiento quedara libre
otra compra podría tomarlo y el \texttt{Commit} llegaría a un asiento ya
vendido.

Después de votar \texttt{Yes} el participante queda en
\texttt{EstadoParticipante::Preparado} y pierde la autonomía. Confirmar no
puede, porque no sabe cómo votó el otro. Abortar tampoco, porque el
coordinador pudo haber decidido \texttt{Commit}. Le queda esperar, y el bucle
del actor lo refleja en el tipo de espera:

\begin{lstlisting}
let msg = if estado == EstadoParticipante::Preparado {
    match rx.recv_timeout(timeout_decision) {
        Ok(m) => m,
        // Estoy PREPARADO: no puedo hacer commit ni abort por mi cuenta.
        // Anoto que quede bloqueado y sigo esperando la decision.
        Err(RecvTimeoutError::Timeout) => { bloqueado = true; continue; }
        Err(RecvTimeoutError::Disconnected) => break,
    }
} else {
    match rx.recv() { Ok(m) => m, Err(_) => break }
};
\end{lstlisting}

El \texttt{continue} hace que al vencer el timeout el participante anote el
bloqueo en su \texttt{wal} y vuelva a esperar. Ese \texttt{wal} es
un vector de líneas que en un sistema real sería un log en disco, el que le
permite a un servicio que se reinicia recordar que había votado \texttt{Yes}
en lugar de abortar por defecto.

\subsection{Fase 1: Prepare}

El coordinador manda \texttt{Prepare} a todos con un canal de respuesta
compartido y junta los votos contra una fecha límite única, no contra un
timeout por participante:

\begin{lstlisting}
// FASE 1 - PREPARE: pregunto a todos y junto los votos con timeout.
let (tx_votos, rx_votos) = channel::<(String, Voto)>();
for (_, p) in participantes {
    let _ = p.send(MsgParticipante::Prepare { tx, responder_a: tx_votos.clone() });
}
let limite = Instant::now() + timeout_votos;
let mut votos: HashMap<String, Voto> = HashMap::new();
while votos.len() < participantes.len() {
    let restante = limite.saturating_duration_since(Instant::now());
    match rx_votos.recv_timeout(restante) {
        Ok((nombre, v)) => { votos.insert(nombre, v); }
        // Vencio el plazo: los que faltan cuentan como NO.
        Err(RecvTimeoutError::Timeout) => break,
        Err(RecvTimeoutError::Disconnected) => break,
    }
}
let todos_yes = participantes.iter()
    .all(|(n, _)| votos.get(n) == Some(&Voto::Yes));
\end{lstlisting}

Al participante que no contesta dentro del plazo se lo cuenta como
\texttt{No}, porque no llega un mensaje de error sino que falta un mensaje,
y esa ausencia es todo lo que el coordinador tiene para interpretar. La
lectura es segura porque quien no contestó no reservó nada, y si llegó a
reservar va a recibir el \texttt{Rollback} de todas formas. Contar esa falta
de respuesta como \texttt{Yes} habilitaría un \texttt{Commit} sobre un
recurso que nadie garantizó.

\subsection{Fase 2: Commit o Rollback}

Con los votos juntos el coordinador decide una sola vez:

\begin{lstlisting}
// FASE 2 - COMMIT / ROLLBACK: la decision es unica y final.
let decision = if todos_yes { Decision::Commit } else { Decision::Rollback };
for (nombre, p) in participantes {
    let msg = match decision {
        Decision::Commit => MsgParticipante::Commit { tx },
        Decision::Rollback => MsgParticipante::Rollback { tx },
    };
    let _ = p.send(msg);
}
\end{lstlisting}

El \texttt{Commit} necesita unanimidad, y un solo \texttt{No} o un solo voto
que no llega alcanzan para el \texttt{Rollback}. La decisión se registra
antes de enviarse. Si el coordinador cayera después de mandar el primer
\texttt{Commit} y al reiniciarse no recordara qué decidió, podría mandarle
\texttt{Rollback} al segundo participante y partir la transacción en dos
mitades inconsistentes.

Del lado del participante, \texttt{Commit} se aplica solo si estaba
\texttt{Preparado} y \texttt{Rollback} se aplica siempre que no haya
confirmado. \texttt{revertir} es idempotente, así que un \texttt{Rollback}
que le llega a quien votó \texttt{No} no hace nada.

\subsection{Fallas inyectadas}

\texttt{ejecutar\_escenario} arma la compra a partir de un
\texttt{Escenario} con el saldo, el asiento y las fallas:
\texttt{Comportamiento::NoResponde} para un participante que deja de
contestar y \texttt{FallaCoordinador::CaeTrasPrepare} para el coordinador
que se cae con los votos en la mano. \texttt{cargo run -- parte3} imprime
cinco casos con el log del coordinador y el \texttt{wal} de cada servicio:
el camino feliz, un \texttt{No} de cada servicio, el participante que no
contesta y la caída del coordinador. Son los mismos cinco que verifican los
tests.

\subsection{Partición de red después del Prepare}

En el caso E los dos servicios votaron \texttt{Yes} y quedaron
\texttt{Preparado}, con los fondos retenidos y el asiento bloqueado. Ahí se
parte la red y la orden de \texttt{Commit} no llega.

Desde donde están, los participantes no saben si el coordinador alcanzó a
decidir \texttt{Commit} o si decidió \texttt{Rollback} por el timeout de
algún otro. Con esa duda no pueden moverse. Si confirman y la decisión era
\texttt{Rollback}, o si abortan y era \texttt{Commit}, la transacción queda
partida en un asiento vendido sin cobrar o un cobro sin asiento.
Consultarse entre ellos no ayuda, porque los dos están en el mismo estado, y
preguntarle al coordinador es lo que la partición impide.

Se quedan esperando, y mientras esperan siguen reteniendo el recurso. El
asiento no se le puede vender a nadie más y la plata del usuario no está ni
cobrada ni disponible. En la salida del caso E eso aparece como
\texttt{fondos\_retenidos=true} y \texttt{asiento\_retenido=true}. Una
transacción trabada le bloquea recursos a las que vienen atrás.

En términos de CAP, 2PC es un protocolo CP, porque ante la partición
conserva la consistencia y pierde disponibilidad. El sistema deja de responder compras sobre esos
recursos antes que arriesgar un asiento vendido dos veces.

\subsection{Safety o liveness}

Se compromete \emph{liveness}; \emph{safety} se mantiene.

Safety acá es la atomicidad: todos los participantes terminan con la misma
decisión y nunca hay uno confirmado y otro abortado. Durante la partición
eso sigue valiendo, y sigue valiendo porque los participantes permanecen a
la espera. El bloqueo es el mecanismo con el que 2PC preserva la atomicidad.
Un participante que al vencer el timeout decidiera por su cuenta estaría
eligiendo disponibilidad y podría dejar la transacción partida.

Liveness es la terminación: toda transacción debería llegar a un estado
final en tiempo finito. Eso se pierde porque los participantes quedan en un
estado incierto y no hay cota para cuánto van a seguir ahí. Esta es la propiedad
\emph{blocking} de 2PC y es inherente al protocolo, porque la información
que el participante necesita para salir de \texttt{Preparado} está
únicamente en el coordinador.

La ventana se puede angostar, no eliminar. Un protocolo cooperativo de
terminación deja que los participantes se consulten: si alguno recibió la
decisión la propaga, y si alguno todavía no votó se puede abortar. No
alcanza cuando todos están \texttt{Preparado}, que es el caso E. Lo que se
usa en la práctica es replicar al coordinador con un algoritmo de consenso,
de manera que la decisión quede en varios nodos y el reemplazo la lea y la
reenvíe. Si se pierde la mayoría el sistema vuelve a bloquearse.

\subsection{Pruebas}

Las pruebas están en el módulo \texttt{dos\_fases} y se corren con
\texttt{cargo test --lib dos\_fases}. Cada una fija un escenario y revisa la
decisión del coordinador y el reporte final de los dos participantes.

\texttt{todos\_yes\_hace\_commit} es el camino feliz. La decisión es
\texttt{Commit}, los dos quedan \texttt{Confirmado}, el saldo es 950 y el
asiento vendido.

\texttt{un\_no\_en\_inventario\_hace\_rollback\_y\_libera\_fondos} arranca
con el asiento ya vendido. La decisión es \texttt{Rollback} y lo que se mira
es la compensación: el saldo volvió a 1000 y Pagos no retiene nada. Un
\texttt{Yes} no deja plata colgada si la transacción aborta.

\texttt{saldo\_insuficiente\_hace\_rollback\_y\_libera\_asiento} es el
simétrico, con un saldo de 10, y verifica que el asiento que Inventario
había bloqueado vuelva a estar libre, porque un asiento que no se libera es
una entrada que no se puede vender.

\texttt{participante\_que\_no\_responde\_provoca\_rollback\_por\_timeout}
pone a Pagos en \texttt{NoResponde}. El voto que falta queda en
\texttt{None}, la decisión es \texttt{Rollback} e Inventario libera el
asiento; el coordinador no se queda esperando un voto que no va a llegar.

\texttt{coordinador\_cae\_tras\_prepare\_bloquea\_a\_los\_participantes}
cubre el caso analizado. Verifica que no hubo decisión (\texttt{None}), que
los dos quedaron \texttt{Preparado} y marcados como bloqueados y que los dos
siguen reteniendo su recurso. Lo que se afirma ahí es que el protocolo no
termina y que ninguno de los dos se adelantó a decidir solo.
